\colorlet{Blue}{FigureBlue}
\colorlet{Green}{FigureGreen}
\colorlet{Red}{FigureRed}
\colorlet{Stroke}{FigureInk}
\colorlet{BlueFill}{FigureBlue!14}
\colorlet{GreenFill}{FigureGreen!14}
\tikzset{every node/.style={font=\figurefont\fontsize{8}{9.5}\selectfont},note/.style={font=\figurefont\fontsize{7.5}{9}\selectfont}}
% Scalar, vector, matrix and order-3 array as index structures; final width 112 mm.
% One highlighted cell per panel shows how many indices locate a single entry.
\begin{tikzpicture}[x=1mm,y=1mm,
  cell/.style={draw=Stroke,fill=white,line width=.6pt},
  hit/.style={draw=Blue,fill=BlueFill,line width=.8pt}]
\path[use as bounding box] (0,-6) rectangle (112,30);
\def\c{4.2}
% ---- scalar
\draw[hit] (8,10) rectangle +(\c,\c);
\node at (8+\c/2,10+\c/2) {$a$};
\node at (10,-3) {标量};
% ---- vector: 4 entries as a column
\foreach \i in {0,...,3}{\draw[cell] (30,{18-\i*\c}) rectangle +(\c,\c);}
\draw[hit] (30,{18-2*\c}) rectangle +(\c,\c);
\node[anchor=west,text=Blue] at (30+\c+1,{18-2*\c+\c/2}) {$x_i$};
\node at (33,-3) {向量};
% ---- matrix: 4 x 5 grid
\foreach \i in {0,...,3}{\foreach \j in {0,...,4}{
  \draw[cell] ({50+\j*\c},{18-\i*\c}) rectangle +(\c,\c);}}
\draw[hit] ({50+3*\c},{18-1*\c}) rectangle +(\c,\c);
\node[anchor=south,text=Blue] at ({50+3.5*\c},{18+\c}) {$A_{i,j}$};
\draw[Blue,line width=.6pt,->] ({50+3.5*\c},{18+\c+.3}) -- ({50+3.5*\c},{18-1*\c+\c+.2});
\node at (60.5,-3) {矩阵};
% ---- order-3 array: three 4 x 4 slices, offset back-to-front
\foreach \k/\dx in {2/6,1/3,0/0}{
  \foreach \i in {0,...,3}{\foreach \j in {0,...,3}{
    \draw[cell] ({84+\dx+\j*\c},{18+\dx*0.8-\i*\c}) rectangle +(\c,\c);}}}
\draw[hit] ({84+1*\c},{18-2*\c}) rectangle +(\c,\c);
\node[anchor=north,text=Blue] at ({84+1.5*\c},{18-3*\c-.4}) {$X_{i,j,k}$};
\node at (95,-3) {三阶数组};
\end{tikzpicture}
